% Standalone section fragment for the accompanying report.
\section{The seven-state quotient: a global finite certificate}

Let $R$ be the explicitly supplied unfolded 202-state table and let $M$
be \texttt{pa195.native.tm}.  Both use native start label $74$ and native
halt label $254$.  The dense presentation of $M$ has start $0$, working
labels $0,\ldots,194$, and halt label $195$.
The quotient map $\sigma$ identifies the following seven pairs and fixes
all other controls.  The first member of each pair is removed.
\[
 (2018,2024),\ (2030,2050),\ (2038,2019),\ (2042,2028),\
 (2053,2043),\ (2055,2034),\ (2062,2036).
\]
The unfolded machine consists of 134 inherited controls and 68 newly
allocated AX reader controls.  All seven identifications are disjoint, so
$202-7=195$ working controls remain.  The root-reachability census visits
every one of the 195 controls; no distinguished halt control is counted.

\paragraph{The incoming-history invariant.}
For every pair $(a,b)$ in this list, every incoming row of $R$ whose target
is $a$ writes $0$ and moves right; every incoming row whose target is $b$
writes $1$ and moves right.  There are exactly 22 such incoming rows, all
listed in \texttt{local\_fold\_certificate.json}.  Since neither $a$ nor
$b$ is the start control, every reachable configuration of $R$ satisfies
\[
 q=a\ \Longrightarrow\ t(h-1)=0,
 \qquad
 q=b\ \Longrightarrow\ t(h-1)=1,
\]
where $h$ is the head position and $t$ is the tape.  This is a direct
one-step induction.  It applies to every initial tape, not only to the
blank run or to well-formed register-program frames.

\paragraph{The 28 local obligations.}
For each of the 14 old controls and each current bit $v\in\{0,1\}$,
start $M$ in control $\sigma(q)$ at head position $0$, with the prescribed
history bit at $-1$ and current bit $v$ at $0$.
The certificate gives a positive finite run that reproduces the single
$R$ transition, including its write, head displacement, and projected
successor.  Until that endpoint every read and write is at $-1$ or $0$.
The two-cell tape at the endpoint is compared exactly.  Hence every
extension of those two cells to a bi-infinite tape satisfies the same
identity.  Of the 28 runs, ten have length 1, seventeen have length 3,
and one has length 5.  All other old transitions commute with $\sigma$
in exactly one step.

\paragraph{Halting preservation in both directions.}
Map a configuration $(q,h,t)$ to $(\sigma(q),h,t)$.  The preceding finite
identities and incoming-history invariant show that every step of $R$
from a reachable configuration is simulated by one to five steps of $M$.
Every intermediate control has an outgoing transition unless the final
endpoint is the halt control; $\sigma$ fixes that halt control.
Concatenation therefore sends any $R$ halting trace to an $M$ halting
trace.

Conversely, suppose $M$ halts after $T$ transitions.  If $R$ does not halt
within its first $T$ transitions, its first $T+1$ steps exist.  Their
simulating blocks concatenate to more than $T$ defined $M$ transitions,
in contradiction with determinism and the absence of outgoing rows at
halt.  Consequently $R$ halts within $T$ steps.  This also covers halts
inside an attempted comparison block; an intermediate premature halt
would contradict its verified finite trace.

All quantified tape and trace operations in this proof have standard
primitive-recursive finite encodings.  The invariant is an induction on
trace length, the local lookup consists of 28 fixed finite cases, the
forward trace construction concatenates blocks of length at most 5, and
the converse uses bounded iteration through $T+1$ parent steps.  Thus the
argument formalizes in PA and proves
\[
 \mathrm{PA}\ \vdash\
 \operatorname{Halt}(R,\mathrm{blank})
 \ \leftrightarrow\
 \operatorname{Halt}(M,\mathrm{blank}).
\]
No assumption of consistency, standard-model soundness, or nonhalting
enters this quotient theorem.
